\chapter[Referencial Teórico]{Referencial Teórico}

\section{MySQL}

De acordo com \citeonline{milani2007}, o cenário de globalização e o avanço da automação, impulsionados pela popularização da internet, tornaram o armazenamento de dados um passo primordial para a migração de negócios para o ambiente digital. Nesse contexto, o MySQL surge para suprir essa demanda de organização de informações.

O MySQL é definido como um Sistema Gerenciador de Banco de Dados (SGBD) relacional de licença dupla, incluindo uma versão de código aberto (open source). Embora tenha sido projetado inicialmente para aplicações de pequeno e médio porte, a ferramenta evoluiu para suportar projetos de grande escala. Segundo \cite{milani2007}, o MySQL é reconhecido por sua robustez, apresentando um nível de \textbf{paridade} técnica (ou nivelamento) com softwares proprietários de mercado, tais como o SQL Server e o Oracle.

\section{PHP}
Conforme afirma \citeonline{DallOglio2015}, o PHP — que originalmente significava Personal Home Page — é uma linguagem de programação criada por Rasmus Lerdorf no outono de 1994. Inicialmente, a ferramenta consistia em um conjunto de scripts escritos em linguagem C, destinados à criação de páginas dinâmicas para que o autor monitorasse o acesso ao seu currículo na internet. Com o aumento da popularidade entre outros usuários, Lerdorf adicionou novos recursos, como a interação com bancos de dados. Em 1995, o código-fonte foi liberado, permitindo a colaboração de novos desenvolvedores. Vale ressaltar que, durante um curto período, a linguagem também foi denominada FI (Forms Interpreter).
\section{Laravel}

\section{VsCode}